\documentclass[exercice,a4]{thedocument}%
\usepackage{texmaths-tikz}%
\usepackage{pstricks}%
\usepackage{pst-bezier}%
\usepackage{pstricks-add}%
\usepackage{pst-eucl}%
\usepackage{pst-func}%
\usepackage{pst-tree}%
\usepackage{pst-circ}%
\usepackage{tikz}%
\usepackage{pgfplots}%
\usepackage{tkz-tab}%
\usepackage{tikzpagenodes}%
\usepackage{tabularray}%
\usepackage{tabularx}%
\usepackage{multirow}%
\usepackage{multicol}%
\usepackage{scratch3}%
\usepackage{enumitem}%
\usepackage[np]{numprint}%
\usepackage{pas-tableur}%
\usepackage{listings}%
\usepackage{graphicx}%
\usepackage{adjustbox}%
\usepackage{needspace}%
\usepackage{xcolor}%
\usepackage{eurosym}%
\usepackage{diagbox}%
\usepackage{colortbl}%
\usepackage{multido}%
\usepackage{icomma}%
\usepackage[french]{babel}%
\startdatakeys
    \source{26GENMATAG1}
    \BO{2026}
    \cycle{4}
    \competencesDuSocle{CA,RA,CO}
    \objectifsApprentissage{C4A1aa,C4A2ad,C4B2cb,C4D2aa,C4D3ac}
    \motsCles{programme de calcul,Scratch,expression littérale,équation}
\enddatakeys
\definecolor{gris}{gray}{0.5}%
\providecommand{\wideparen}[1]{\overset{\frown}{#1}}%
\begin{document}%
\providecommand{\degres}{^\circ}%
On considère le programme de calcul ci-dessous.

\begin{center}
	\fbox{
		\begin{minipage}{0.55\linewidth}
			\centering \textbf{Programme A}
			\begin{itemize}[label=$\bullet~$]
				\item Choisir un nombre
				\item Soustraire 2
				\item Multiplier par 5
				\item Ajouter le triple du nombre de départ
			\end{itemize}
		\end{minipage}
	}
\end{center}

\begin{enumerate}
	\item Vérifier qu'en choisissant 1 comme nombre de départ, on obtient $-2$ avec ce programme de calcul.

	\item On note $x$ le nombre choisi au départ.
	\begin{enumerate}
		\item Montrer que le résultat du programme est $8x - 10$.

		\item Quel nombre de départ faut-il choisir pour que le résultat du programme soit égal à 0 ?
	\end{enumerate}

	\item \emph{Question d'algorithmique}
	\begin{enumerate}
		\item On considère ci-dessous le programme de calcul B, réalisé avec le logiciel \og Scratch \fg.

		\begin{center}
			%\fbox{
				%\begin{minipage}{0.6\linewidth}
				\centering \textbf{Programme B}\\[6pt]

				\begin{scratch}
					\blockinit{Quand \greenflag est cliqué}
					\blocksensing{demander \ovalnum{choisir un nombre} et attendre}
					\blockvariable{mettre \selectmenu{résultat} à \ovaloperator{\ovalnum{4}*\ovalvariable{réponse}}}
					\blockvariable{mettre \selectmenu{résultat} à \ovaloperator{\ovalvariable{résultat} - \ovalnum{5}}}
					\blockvariable{mettre \selectmenu{résultat} à \ovaloperator{\ovalnum{2} *\ovalvariable{résultat}}}
					\blocklook{dire \ovalvariable{résultat}}
				\end{scratch}
				%\end{minipage}
				%}
		\end{center}

		En appliquant ce programme de calcul au nombre 1, quel résultat final obtient-on ?

		\item Montrer que le résultat du programme A est toujours le même que celui du programme B quel que soit le nombre choisi au départ.
	\end{enumerate}
\end{enumerate}

\needspace{10\baselineskip}
\medskip
\showthesource
\end{document}
